\clearpage
\begin{explanationbox}
\begin{center}
 \begin{tikzpicture}[x=1cm,y=1cm,>=stealth,font=\small]
\path[use as bounding box] (-5.4,-3.8) rectangle (5.4,4.1);
\node[font=\bfseries] at (0,3.8) {中心对称：两个斜率交换，乘积不变};
\draw[->,black!50] (-4.6,0)--(4.6,0) node[right] {$x$};
\draw[->,black!50] (0,-3.1)--(0,3.2) node[above] {$y$};
\foreach \s in {-1,1} {\draw[blue!65!black,thick,domain=-3:3,samples=80,variable=\t] plot ({\s*sqrt(4+.8*\t*\t)},\t);}
\draw[orange!85!black,thick] (-2,0)--(3,2.5) (2,0)--(-3,-2.5);
\draw[teal!80!black,thick] (2,0)--(3,2.5) (-2,0)--(-3,-2.5);
\fill (-2,0) circle (2pt) node[above left] {$A_1$};
\fill (2,0) circle (2pt) node[below right] {$A_2$};
\fill (3,2.5) circle (2pt) node[right=8pt] {$P_1$};
\fill (-3,-2.5) circle (2pt) node[left=8pt] {$P_2$};
\node[orange!85!black,fill=BGExplanation] at (-.9,1.35) {$1/2$};
\node[teal!80!black,fill=BGExplanation] at (3.6,1.25) {$5/2$};
\node[teal!80!black,fill=BGExplanation] at (-3.6,-1.25) {$5/2$};
\node[orange!85!black,fill=BGExplanation] at (.9,-1.35) {$1/2$};
\node at (0,-3.6) {$P_1=(3,5/2),\quad P_2=(-3,-5/2),\quad k_1k_2=5/4$};
\end{tikzpicture}

 
 {\small\color{black!65}示意：$\frac{x^2}{4}-\frac{y^2}{5}=1$，
 $P_1(3,\frac52)$、$P_2(-3,-\frac52)$ 均在曲线上；实线是对应顶点连线。}
\end{center}
\textbf{\textcolor{CExplanation}{一句话直觉}}
两条连线的斜率有\textbf{相同的竖直高度} \(y\)，
分母分别是 \(x+a\) 与 \(x-a\)；相乘后自然出现\textbf{平方差} \(x^2-a^2\)。
双曲线方程恰好说明 \(y^2\) 就是这一平方差的固定倍数。

\medskip
\textbf{符号观察：}双曲线非顶点处 $|x|>a$，$x+a$ 与 $x-a$ 同号，
两斜率同号，故乘积为正。图中中心对称的两点交换斜率值，乘积仍为 $5/4$。

\textbf{\textcolor{CExplanation}{最省心的代数连接}}
\[
 \underbrace{k_{PA_1}k_{PA_2}}_{\text{几何：两条斜率相乘}}
 =\frac{y^2}{(x+a)(x-a)}
 =\frac{y^2}{x^2-a^2}
 \overset{P\in C}{=}\frac{b^2}{a^2}.
\]
\textbf{识别信号：}题中同时出现“\textbf{实轴两端点}\(+\)\textbf{曲线上一点}\(+\)\textbf{两条连线的斜率}”，
优先想第三定义，不要急于设一般直线与双曲线联立。

\medskip
\textbf{\textcolor{CExplanation}{横纵轴型统一理解（比死记更稳）}}
\begin{itemize}[leftmargin=*,itemsep=0.35em]
 \item 横轴型 \(\frac{x^2}{a^2}-\frac{y^2}{b^2}=1\)：
 斜率 \(=\Delta y/\Delta x\)，故乘积 \(=b^2/a^2=e^2-1\)。
 \item 纵轴型 \(\frac{y^2}{a^2}-\frac{x^2}{b^2}=1\)：
 斜率 \(=\Delta y/\Delta x\)，故乘积 \(=a^2/b^2=1/(e^2-1)\)。
 \item \textbf{不用记横纵对调的结果：}先写两条斜率，再用标准方程替换平方差，只需两步。
\end{itemize}

\textbf{\textcolor{CExplanation}{与椭圆第三定义的联系}}
椭圆 \(\frac{x^2}{a^2}+\frac{y^2}{b^2}=1\)（\(a>b>0\)）
关于实轴端点的斜率积为 \(-b^2/a^2\)（非端点处），
而双曲线为 \(+b^2/a^2\)。\textbf{关键不是背正负号，而是看}
\(y^2\) 与 \(x^2-a^2\) 在曲线上是否同号。

\textbf{\textcolor{CExplanation}{逆向思考的价值}}
已知\(k_1k_2=m>0\)，就得到
\(y^2=m(x^2-a^2)\)；这不只是用来计算定值，
还可\textbf{直接建立点的轨迹方程}。

\medskip
\end{explanationbox}
\clearpage
\begin{explanationbox}
\textbf{\textcolor{CExplanation}{从初中圆幂定理走到双曲线（新增几何桥梁）}}
\begin{center}% Exact model: a=2, b=sqrt(5), P=(3,5/2), H=(3,0).
% T=(4/3,-2sqrt(5)/3), so OT^2=4, OT dot HT=0, HT^2=5.
\begin{tikzpicture}[x=1.1cm,y=1.1cm,>=stealth,font=\small]
\draw[rounded corners=5pt,black!25] (-3.7,-3.65) rectangle (5.1,4.5);
\node[font=\bfseries] at (.7,3.9) {圆幂桥梁：共同高 $PH$ 与切线长 $HT$};
\coordinate (O) at (0,0);
\coordinate (A1) at (-2,0); \coordinate (A2) at (2,0);
\coordinate (H) at (3,0); \coordinate (P) at (3,2.5);
\coordinate (T) at ({4/3},{-2*sqrt(5)/3});
\draw[->,black!50] (-3.15,0)--(4.55,0) node[right] {$x$};
\draw[->,black!50] (0,-2.25)--(0,3.25) node[above] {$y$};
\draw[blue!65!black,thick,domain=0:3.15,samples=80,variable=\t]
 plot ({sqrt(4+.8*\t*\t)},\t);
\draw[blue!65!black,thick,domain=0:3.15,samples=80,variable=\t]
 plot ({-sqrt(4+.8*\t*\t)},\t);
\draw[violet!75!black,dashed,thick,domain=0:360,samples=120,variable=\t]
 plot ({2*cos(\t)},{2*sin(\t)});
\draw[orange!85!black,thick] (A1)--(P);
\draw[teal!80!black,thick] (A2)--(P);
\draw[black!75,very thick] (P)--(H);
\node[right=7pt] at (3,1.25) {$PH=|y|$};
\draw[red!65!black,very thick] (H)--(T);
\node[red!65!black] at (2.45,-1.25) {$HT$};
\draw[violet!75!black,thick] (O)--(T);
\node[below left=2pt] at (.65,-.72) {$a$};
\draw[black!60] ($(T)!.18cm!(O)$)--($(T)!.18cm!(O)+(T)!.18cm!(H)-(T)$)--($(T)!.18cm!(H)$);
\draw (2.82,0)--(2.82,.18)--(3,.18);
\fill (O) circle (1.6pt) node[above left=3pt] {$O$};
\fill (A1) circle (1.8pt) node[below left=5pt] {$A_1$};
\fill (A2) circle (1.8pt) node[below left=5pt] {$A_2$};
\fill (H) circle (1.8pt) node[below right=5pt] {$H(x,0)$};
\fill (P) circle (1.8pt) node[right=8pt] {$P(x,y)$};
\fill (T) circle (1.8pt) node[below right=7pt] {$T$};
\node[violet!75!black] at (-2.75,-1.7) {辅助圆};
\draw[dotted,black!35] (-2,-.3)--(-2,-2.6) (3,-.4)--(3,-3.1) (2,-.4)--(2,-3.1);
\draw[<->,black!70] (-2,-2.6)--(3,-2.6) node[midway,fill=white,inner sep=3pt] {$HA_1=x+a$};
\draw[<->,black!70] (2,-3.1)--(3,-3.1);
\node[left=5pt] at (2,-3.1) {$HA_2=x-a$};
\end{tikzpicture}
\end{center}
图取右支 $P$，$H$ 为垂足；辅助圆以 $A_1A_2$ 为直径，半径为 $a$。
$|x|>a$ 保证 $H$ 在圆外；图中 $OT\perp HT$，故 $HT$ 是切线。
\textbf{初中切割线定理}给出
\[
 HA_1\cdot HA_2=HT^2=OH^2-a^2=x^2-a^2.
\]
两个直角三角形的\textbf{共同高} $PH=|y|$，所以
\[
 \underbrace{k_{PA_1}k_{PA_2}}_{\text{两斜率同号}}
 =\frac{PH^2}{HA_1\,HA_2}
 =\boxed{\left(\frac{PH}{HT}\right)^2}.
\]
而双曲线方程恰好是 $PH^2=(b^2/a^2)HT^2$。
因此\textbf{斜率积不再只是代数约分，而是“共同高与切线长的比值平方”}。
\textbf{易懂但不冒充纯几何：}圆幂步骤是几何的，$PH/HT=b/a$ 仍由双曲线标准方程确定。

\medskip
\end{explanationbox}
\clearpage
\begin{explanationbox}
\textbf{\textcolor{CExplanation}{一个斜率积，统一圆、椭圆和双曲线}}
% [TikZ待绘] C065_unified_locus.tex：相同A1,A2，m=-1圆、-1/4椭圆、5/4双曲线。
\IfFileExists{assets/tikz/C065_unified_locus.tex}{%
 \begin{center}\input{assets/tikz/C065_unified_locus.tex}\end{center}
}{}
对\textbf{任意}动点 $P(x,y)$（仅要求 $x\ne\pm a$），不预设其曲线类型：
\[
 k_1k_2=m\iff y^2=m(x^2-a^2).
\]
\begin{itemize}[leftmargin=*,itemsep=.22em]
 \item $m<0$：$\dfrac{x^2}{a^2}+\dfrac{y^2}{a^2|m|}=1$，轨迹是\textbf{椭圆去掉两个端点}；
 特别地 $m=-1$ 时，$x^2+y^2=a^2$ 为\textbf{圆（同样去端点）}。
 \item $m>0$：$\dfrac{x^2}{a^2}-\dfrac{y^2}{a^2m}=1$，轨迹是\textbf{双曲线去端点}。
 \item $m=0$：$y=0$ 且 $x\ne\pm a$，是\textbf{$x$ 轴去掉两个端点}。
\end{itemize}
圆的情况立即联系初中\textbf{直径所对圆周角是直角}：
\[
 k_1k_2=-1\iff PA_1\perp PA_2
 \iff P\text{ 在以 }A_1A_2\text{ 为直径的圆上（非端点）}.
\]
由此理解 C064 椭圆斜率积为负、C065 双曲线斜率积为正。
注意：当 $|m|>1$ 且 $m<0$，椭圆长轴在 $y$ 方向；$A_1,A_2$ 只是 $x$ 轴交点。

\medskip
\textbf{\textcolor{CExplanation}{半角视角和渐近线：为什么会趋近但取不到等号}}
右支作三角参数 $x=a\sec\theta,\ y=b\tan\theta$
（$-\pi/2<\theta<\pi/2$，$\theta\ne0$），
则两个斜率正好是
\[
 k_1=\frac ba\tan\frac\theta2,\qquad
 k_2=\frac ba\cot\frac\theta2.
\]
\textbf{互为倒数的半角因子}使积直接成为 $b^2/a^2$。
同时当 $|x|\to\infty$ 时，两斜率都趋向同一条渐近线的斜率
$\pm b/a$，但在有限点处永远不相等。
因而若 $s=k_1+k_2$、$m=k_1k_2>0$，则
\[
 s^2-4m=(k_1-k_2)^2>0,
\]
即 $|s|>2\sqrt m$；边界 $|s|=2\sqrt m$ \textbf{只能逼近，不能达到}。

\end{explanationbox}
